\section{The Image and Likeness}

\begin{quotex}
When this thought or contemplation visits the inner temple of our spirit, this is filled with the light and majesty of God. The spirit's rapture then is indescribable. All its bones, that is, all the smallest moving parts of its being, give an irrepressible leap of deep inward joy. \flright{\textsc{Bishop Theophan the Recluse}}

Blessed are the pure in heart, for they shall see God. \flright{\textit{Mt 5:8}}

\end{quotex}
It is said that man is created in the image and likeness of God. Although the West tends to conflate the two terms, that is not quite correct esoterically. In what follows, please refer to the schema~\ref{tab:ImageLikeness}.

\begin{table}[h]\centering\footnotesize
\label{tab:ImageLikeness}
\begin{tabular}{rcccc}\toprule
 &
\textbf{Exoteric Teaching} &
\textbf{Esoteric Teaching} &
\textbf{Source} &
\textbf{Reality}\\\midrule
\textbf{Image} &
Intellect and Will &
Real I and True Will &
Intrinsic &
Virtual\\
\textbf{Likeness} &
Grace &
Transcendent influences &
Acquired &
Actual\\\bottomrule
\end{tabular}
\caption{Image and likeness}
\end{table}
The image of God can only be perfected in a purified consciousness, freed of lower influences, by the action of Grace. ``Pop" religion today regards the ``image and likeness" as a fait accompli, asserting a certain dignity on those. However, the image of God in man is only virtual and needs to be actualized. That is the ``likeness" of God, the process by which man becomes more god-like, or theosis. He becomes a ``son of God", which means sharing in the same nature. To become such a son by adoption is a privilege and is not at all universal. Pop religion, on the other hand, claims that everyone is a ``child of god" no matter what. The opposite misconception is that God is ``totally other" to man. While that may be true for profane man, it is not so for those in communion with God.

Now, following the Greek Fathers, the Eastern Church\footnote{\url{http://www.stgeorgeserbian.us/darren/darren03.htm}} has been clearer about making this distinction. The failure to take this into consideration often leads to mischief in social policies.

\paragraph{The Image of God}
So, to summarize, the ``image" is what we are born with, but it is as yet only virtual. As the image of God, man has intelligence and free will. Of course, that is the higher intellect, not just discursive thought. This higher intellect is called the ``heart" in Tradition. When the heart is purified, that is, purged of the distorting effects of negative emotions, disordered desires, and false doctrines, then the image of God will be more perfectly reflected in one's consciousness. In that state, one is ``pure of heart".

\paragraph{The Likeness of God}
As one overcomes the lure of the lower influences, one becomes open to ``grace", or transcendent influences. One's nature changes to become more God-like. One is given the gift of understanding. One's will becomes free, since it is not being driven by external forces. One hungers for impartial justice and so on. One becomes passionless and free of useless sentiment. Nevertheless, there can be the experience of great joy, as Theophan describes, and even flights of ecstasy.

Pop religion presumes we are by birth children of God. Misleadingly, they take this from John as an indication:

\begin{quotex}
That was the true light, which enlightens every man that comes into this world. (John 1:9) 

\end{quotex}
As such, it appears universal, however that is immediately followed by:

\begin{quotex}
He was in the world, and the world was made by him, and the world knew him not. He came unto his own, and his own received him not. But as many as received him, he gave them power to be made the sons of God, to them that believe in his name. Who are born, not of blood, nor of the will of the flesh, nor of the will of man, but of God. (John 1:10-13) 

\end{quotex}
But that is not the case; we see only the ``second birth" makes one ``born of God" which can only come from the acceptance of the ``true light". In Tradition, to be ``born of God" means to share the same nature, just as the son shares the nature of the father, i.e., he is a ``chip off the old block". The pop religion, on the other hand, interprets the phrase as referring to perpetual infantilism.

\paragraph{Hermetic Teaching}
\textbf{Valentin Tomberg} in his \emph{Meditations} teaches that the image and likeness of God refers to the human personality. He contrasts that with Oriental systems which lead to a depersonalization, which he rejects. Tomberg is following the Eastern Church in recognizing that growing in the image and likeness of God is actually a process of humanization, of becoming ``truly human", a fuller ``realization of true human nature". In effect, he has ``more being" than he did before.

This is an alchemical process. As such, it transforms the base human into the image and likeness of God through ``sublimation". In Hermetic initiation, one penetrates into the depths of consciousness to awaken to the primordial Adamic state of the image and likeness.

While the contemporary \emph{Catechism of the Catholic Church} shows the loss of any understanding of the distinction between the image and likeness, that was not also the case of the Nordic-Roman Medieval Tradition, which, like the East, had not forgotten that doctrine. Quoting \textbf{St Bernard}, Tomberg shows how the image and likeness separated after the Fall of Adam and the forgetting of the Primordial State. Bernard claims that while the image of God remains in man, the likeness has been lost ``in the soul which sins". As he wrote in Sermon on the Annunciation of the Blessed Virgin Mary:

\begin{quotex}
Man was made in the image and likeness of God: in image he possesses freedom of will, and in likeness he possesses virtues. The likeness has been destroyed; however, man conserves the image. The image can be burnt in hell, but not consumed. It is damaged but not destroyed. Through fate as such it is not effaced, but subsists. Wherever the soul is, there also will be the image. It is not so with the likeness. This remains in the soul which accomplishes the good; in the soul which sins it is wretchedly transformed. The soul which has sinned ranks with beasts devoid of intelligence. 

\end{quotex}
The consequence is that man is ``free", perpetually free through eternity. The likeness, however, is not essentially immortal. The Holy Guardian Angel acts to preserve the likeness that has been lost in man.

\paragraph{Misconceptions}
Distortions arise because of the lack of understanding about the distinction between the image and likeness. In particular, new age and liberal theory more or less adopt the notions of the ``image and likeness" as well as the universalization of the ``children of god". For example, if the image of God is actual rather than virtual at birth, then it would follow that one's weakness, faults, and foibles are in reality God-given attributes of one's personality or nature. This is a reversion to pantheistic paganism as we showed in \textit{Esoteric Stoicism}\footnote{\url{https://gornahoor.net/?p=8147}}. Specifically, one's ``presuppositions" are seen as God given, while in the Christian view such presuppositions need to be examined and tested.

\paragraph{Social Policy}
You will be hearing in the coming month about the ``option for the poor". It is said that there is an inherent ``dignity of man" (i.e., worthy of respect) due to his being in the ``image and likeness of God". Of course, we treat everyone with the respect due him as a matter of justice. Even if a man has lost the likeness of God, ``every sinner has a future" (Augustine). Similarly, the unborn is also due justice, even if his being is still virtual.

In economic policies, there is a religious left that makes common cause with the secular left. On a certain level, this is quite strange. The religious left bases its program on man as being in the image of God and a descendant of God. The secular left, on the other hand, regards man as a material being and a descendant of monkeys. When two incompatible premises lead to the same conclusion, then, by Occam's razor, one is justified in accepting the simpler premise, which is the secular worldview. Hence, the religious left tends over time to get sucked into the secular, material view.

Now, economic solutions are not as simple as many want to believe. There are several alternatives, not necessarily incompatible with each other, to provide for the material needs of the poor. Moreover, the option for the poor should include cultural and spiritual components. This is where the secular worldview falls short, since only economic and material considerations matter. And the religious left, but accepting de facto the secular premise, becomes impotent.

By failing to recognize that the loss of the ``likeness of God" has been lost (due to faulty understanding), the religious left has no solution to the cultural and spiritual issues. The secular left promotes contempt for religion, sexual license, drug use, gun violence, and so on. While in their personal lives, they are usually less affected by the consequences of that worldview, that is not true of the poorer elements which have fewer options. The religious left seems to say nothing about that.

\paragraph{Postscript}
This is not to say that the religious right has much better to offer. By glorifying capitalism and perpetual war, they create more problems than they solve. Even if they are marginally better on cultural and spiritual issues, there is a disturbingly mixed message.

At a recent protest of a Planned Parenthood clinic, the Evangelicals were yelling through a megaphone, presumably to harass the women. I don't know if they checked first to determine if a woman was there due to incest or rape. In another section the Catholics gathered to say the Rosary. The Evangelicals then turned the megaphone onto the Catholics admonishing them ``not to worship Mary." Is there any wonder why the educated left regards the right as unhinged?



\flrightit{Posted on 2015-08-31 by Cologero }

\begin{center}* * *\end{center}

\begin{footnotesize}\begin{sffamily}



\texttt{Br. Giles Mary on 2015-09-07 at 16:26 said: }

As a quasi-celibate (I say quasi because our community is normal in the Church but not yet a canonical ``religious institute" not as if I only practice celibacy occasionally. Gospel purity of heart is of course another matter altogether…) I would like to add my two cents… The Christian celibate, at least, is not consciously transforming raw sexual energy (at least none that I know of) like a practitioner of tantra (I don't think you implied as much, but I only say it for context.) The goal is rather a transformation of vision. 

Christopher West has made a name for himself in the mainstream by popularizing John Paul II's ``Theology of the Body" catecheses. He well makes the point that consecrated celibacy is ideally not like caging a lion. In fact, perhaps things are different in your country, but in America puritanical vision of sex has deep roots and is a major cause of much of our sexual distortion. It seems to be the mainstream wisdom of the Church that the attempt to cage the lion, as you said, is almost certainly doomed to fail. I would like to say a lot more, but I'm short on time. Suffice to say that a full-fledged life decision for a celibate form of life is essential to bearing the proper fruit. Aquinas explains this well in the Summa Theologica. Also, speaking from my own experience, as a sort of quasi-celibate I don't engage in any sexual activity and try to diminish and avoid impure phantasies, but my being is more sexual than ever. Although we're not given in marriage in the Kingdom, nonetheless we are created male and female. There are relevant metaphysical topics regarding nature and grace that would also be interesting…

Also, to the broader thread… We pray at abortion clinics and try to counsel women. I've had some of the most profound and thoroughly dry-unspectacular spiritual experiences there. There is a definite difference between the typically protestant and catholic approaches. It probably has to do, on the one hand, with protestant notions of imputed grace, in contrast to a catholic renovation of the soul. Many can't see beyond their troubled consciences and feel the need to draw some sharp moral lines. It's a self-serving tendency…


\hfill

\texttt{Alistair Fraser on 2015-09-07 at 17:46 said: }

just wanted to tidy up my last paragraph from earlier …………..which leads to the shortening of any romance to the physical by the so called sex addict, they crave the physical orgasmic sensation in conjunction with this discombubulated aftermath reset of thoughts ,which sort of short cuts them out of their depressive cycle and then they get to do it all over again ……..but with the descent into addiction, the thoughts lessen and stagnate, the depression increases and the physical sensation lessens , and just like the gambler and cocaine addict , they keep believing , next time it will be different , Wake up time or deterioration into the abyss.

BR. GILES MARY – your input is valued in extending the understandings , , the puritanical is also over this side, but maybe not so fanatical as over your way in todays times .

you stated ``but my being is more sexual than ever." which is an interesting expression in the celibate status , do you feel that you contain a certain ``contained" vitality or some extra energy as a result of the celibacy and cerebral vision you are commited too ?


\hfill


\hfill

\texttt{Br. Giles Mary on 2015-09-14 at 16:07 said: }

Alistair,

" do you feel that you contain a certain ``contained" vitality or some extra energy as a result of the celibacy and cerebral vision you are commited too ?" 

It's something like this. Certain potentialities of the soul are indeed freed up or activated in this form of life, for sure. An important point that I have to remind even myself of is that we can go too far in conceiving these things as extrinsic to our soul/body composite. 

``We are our bodies…" is, I believe, one of the phases Saint John Paul II drew from the tradition for his Theology of the Body catecheses. He also uses some phrase about our basic psychosomatic makeup in Love and Responsibility. This is a typically Aristotelian notion rather than Platonic, which the Church has pretty well invested herself into. The soul is the form of the body rather than something extrinsic to it. 

So, we even speak of spiritual progress in terms of growth into an integral and perfect human being. Perhaps, at first, it sounds quaint and modern. But, having grown up as a thoroughly truncated modern man, I'm often discovering anew what this means. So, if we're created male and female, and not destined to leave that behind, I think it follows that we express our sexual differentiation with greater depth and clarity the closer we get to beatitude. Besides, who in their right mind would want to become less a man? It's all about Christ and his Bride. Maybe in response, somebody wants to evoke the Adamic-androgyne or something… Anywho…

Glory to God and Thanks for the discussion.


\end{sffamily}\end{footnotesize}
